\documentclass{report}
\usepackage{amsfonts, amsmath}
\newcommand\R{{\mathbb R}}
\newcommand\Q{{\mathbb Q}}
\newcommand\N{{\mathbb N}}
\newcommand\ve{\varepsilon}
\newcommand\Co{{\mathcal C}}
\pagestyle{empty}
\begin{document}
\large
\centerline{\Huge Fisica Matematica II}
\renewcommand{\thefootnote}{ } 
\footnote{\sl Avete 2 ore di tempo. Ogni esercizio
vale dieci punti. La sufficienza si ottiene con un punteggio $\geq$
18. Solo le risposte chiaramente giustificate saranno prese in
considerazione.}
\renewcommand{\thefootnote}{\arabic{footnote}} 
\vskip.1cm
\centerline{\Large Esonero 09-06-2003}
\vskip1cm
\begin{enumerate}
\item Si determini la soluzione $u(x,t)$ di
\[
\begin{split}
&u_{xx}-u=u_t\quad x\in(0,1);\; t>0\\
&u(x,0)=1\quad x\in(0,1)\\
&u(1,t)=e^{-t};\; u_x(0,t)=e^{-t}\quad t>0.
\end{split}
\]
\item   Si determini la soluzione $u(x,t)$ di
\[
\begin{split}
&u_{tt}=u_{xx}+x\quad x\in(0,1);\;t>0\\
&u(x,0)=0;\; u_x(x,0)=0\quad x\in(0,1)\\
&u(0,t)=0;\; u(1,t)=0\quad \;t>0.
\end{split}
\]
\item  Date le funzioni $f_\alpha(x)=\sqrt{\frac \pi 2} e^{-|x|\alpha}$,
$\alpha>0$,  siano
$\hat f_\alpha$ le loro trasformate di Fourier. Si calcoli
\[
\lim_{\alpha\to 0}\int_\R \hat f_\alpha(\xi) g(\xi) d\xi
\]
per ogni $g\in L^\infty(\R)\cap \Co^0(\R)$. 
\item Si consideri l'equazione
\[
\begin{split}
&u_{tt}-u_{xx}=0 \quad x\in\R;\;t>0\\
&u(x,0)=g(x);\; u_x(x,0)=0,
\end{split}
\]
con supp$\,g\subset [-1,1]$. Si definisca l'energia $e(t)$ contenuta
nell'intervallo $[-1,1]$ al tempo $t$
\[
e(t):=\int_{-1}^1 u_t(y,t)^2+u_x(y,t)^2 \;dy.
\]
Si dimostri che $e(t)\leq e(0)$ per ogni $t\geq 0$.
\end{enumerate}
\end{document}
